\begin{table}[!htbp]
\centering
\scriptsize
\setlength{\tabcolsep}{3.5pt}
\resizebox{\textwidth}{!}{%
\begin{tabular}{@{}>{\raggedright\arraybackslash}p{0.31\textwidth}rrrr@{}}
\toprule
Model & \shortstack{Fit raw crossings\\rate (count/\(N\))} & \shortstack{Fit adjustment\\mean/max} & \shortstack{Forecast raw crossings\\rate (count/\(N\))} & \shortstack{Forecast adjustment\\mean/max} \\
\midrule
Joint Q--DESN \(\AL\)--\(\RHS\) & 0.09\% (21/24,000) & 0.0000/0.0813 & 2.52\% (1,197/47,520) & 0.0118/0.6949 \\
Independent Q--DESN \(\AL\)--\(\RHS\) & 0.47\% (114/24,000) & 0.0003/0.1845 & 7.41\% (3,522/47,520) & 0.0166/0.7034 \\
Joint exQDESN \(\exAL\)--\(\RHS\) & 0.00\% (0/24,000) & 0.0000/0.0000 & 0.00\% (0/47,520) & 0.0000/0.0000 \\
Independent exQDESN \(\exAL\)--\(\RHS\) & 0.00\% (0/24,000) & 0.0000/0.0000 & 0.57\% (269/47,520) & 0.0005/0.1887 \\
\bottomrule
\end{tabular}
}%
\caption{Frequency and magnitude of monotone correction for the posterior-mean quantile grids, aggregated over eight simulation settings. Crossing rates are the proportions of adjacent-level comparisons that cross before rearrangement; counts and total opportunities \(N\) are shown in parentheses. Adjustment entries give the mean and maximum absolute change on the simulation response scale induced by the pre-specified monotone rule. Thus the rates measure how often crossings occur, whereas the adjustment summaries describe their magnitude. All rearranged grids have zero crossings.}
\label{qdesn:tab:supp-joint-qdesn-corrected-v4-crossings}
\end{table}
